\documentclass[10pt,a4paper]{amsart}
\usepackage[margin=19mm]{geometry}
\usepackage{amsmath,amssymb,mathtools}
\usepackage[scr=rsfs,cal=euler]{mathalfa}
\usepackage{xfrac}
\usepackage[unicode,hidelinks,bookmarks=false]{hyperref}
\newcommand{\ie}{i.e.\ }
\newcommand{\cf}{cf.\ }
\newcommand{\resp}{respectively\ }
\newcommand{\Subsection}{\subsection}
\providecommand{\nobreakdash}{\nobreak\hbox{-}}
\setlength{\emergencystretch}{2em}
\multlinegap0pt
\usepackage{mathtools}
\usepackage{bm}
\usepackage{amssymb}
\usepackage{enumerate}
\newtheorem{lemma}{Lemma}
\newtheorem{definition}{Definition}
\newtheorem{claim}{Claim}
\title{Inversion diameter and treewidth}
\author{Yichen Wang, Haozhe Wang, Yuxuan Yang, Mei Lu}
\date{Source: arXiv:2407.15384v4 (CC BY 4.0)}
\begin{document}
\maketitle

\noindent\emph{Source and licence.} Inversion diameter and treewidth. Version arXiv:2407.15384v4 (CC BY 4.0), distributed by arXiv under the Creative Commons Attribution 4.0 International licence, http://creativecommons.org/licenses/by/4.0/, as recorded in the arXiv licence metadata for this exact version and shown on the abstract page https://arxiv.org/abs/2407.15384v4. Full licence notice: \texttt{licenses/arxiv-2407.15384-CC-BY-4.0.txt}.\par\smallskip

\noindent\textit{Editorial scope.} Counted unit: the article's lemma independence (source0.tex lines 388--391). The article's complete original proof of it is delivered with it (source0.tex lines 393--494, one \texttt{proof} environment of 102 lines). No mathematics is written, repaired or re-derived: the statement and the proof are the article's own lines, in the article's own notation, and no retained line is substituted or re-flowed.  Retained uncounted as the source-local context the proof needs: lines 292--295; lines 310--315; lines 378--381.  Nothing was omitted from any retained span, so the package carries no omission record.  No retained line contains a citation, so the wrapper carries no bibliography and no bibliography entry.  No retained line contains a figure environment, an image-inclusion command or drawing code, so no image is delivered and the package declares no asset dependency.  Editorial formatting only, in addition: an AMS \texttt{amsart} preamble that restates the article's own declarations and macros used by the retained lines, namely the environments \texttt{lemma}, \texttt{definition}, \texttt{claim} on separate counters, together with the standard packages those lines need.  The retained environments print in this document as Lemma 1, Definition 1, Claim 1 in order of appearance, whereas the article prints the same items with the numbering of its own full document.  The complete native source of the article as distributed in the arXiv e-print for this exact version is bundled unchanged as \texttt{sources/162657/source0.tex}.
\begin{lemma}\label{lemma: clique independence}
    % Suppose $\lambda^{(k)} \le t \le 2k$.
    If there is a $k$-clique of level $m$ with vertices $v_1, \ldots, v_k$ in $G^{(k)}_{m+1}$, then $\mathbf{v}_1, \ldots, \mathbf{v}_k$ are linearly independent.
\end{lemma}

\begin{lemma}\label{lemma: independent or sum}
    % Suppose $\lambda^{(k)} \le t \le 2k$.
    If there is a $k$-clique  in $G^{(k)}_{m+2}$ of level $m$ with vertices $v_1, \ldots, v_k$, and $u$ is a vertex of level $m+1$ adjacent to all ${(v_i)}_{1 \le  i \le k}$,
    then either $\mathbf{v}_1, \ldots, \mathbf{v}_k, \mathbf{u}$ are linearly independent, or $\mathbf{u} = \sum_{i=1}^{k}\mathbf{v}_i$.
    % for every $1 \le j \le k$, $\mathbf{v_1}, \ldots, \mathbf{v_{j-1}}, \mathbf{v_{j+1}}, \ldots, \mathbf{v_k}, \mathbf{u}$ is linearly independent.
\end{lemma}

\begin{definition}\label{D2}
    Let $C$ be a $p$-clique of $G^{(k)}_{m}$ for some $m$. $C$ is called a \emph{bad} clique if $\dim(\textbf{V}_C \cap \textbf{V}_C^{\perp})\ge p-1$, where
    $\textbf{V}_C=\mathcal{L}(\{\mathbf{v}\mid v\in C\})$ and $p\ge 1$.
    \end{definition}

\begin{lemma}\label{lemma: 0 sol independence}
        If there exists a bad $p$-clique of level $m$ in $G^{(k)}_{m+k+2}$ with $p < k$,
    then there exists a bad clique in $G^{(k)}_{m+k+2}$ of size at least $p+1$.
\end{lemma}

\begin{proof}
    We prove it by contradiction. Suppose the $p$-clique $C_1$ with vertices $v_1,v_2,\ldots, v_p$ of level $m$ is the largest bad clique in $G^{(k)}_{m+k+2}$, where $p<k$.  Then $\dim(\textbf{V}_{C_1})=p$ by Lemma~\ref{lemma: clique independence}.
    Let $\textbf{U}=\textbf{V}_{C_1}\cap \textbf{V}_{C_1}^{\perp}$. Then  $\dim(\textbf{U})\ge p-1$ by Definition~\ref{D2}. For every $\mathbf{u}_1, \mathbf{u}_2\in \textbf{U}$, we have $\mathbf{u}_1\perp \mathbf{u}_2$ which means $\textbf{U}$ is self-orthogonal.

    We first show that $\dim(\textbf{U})= p-1$. 
    Suppose otherwise $\dim(\textbf{U})= p$.
    Then $\textbf{U}=\textbf{V}_{C_1} \subseteq \textbf{V}_{C_1}^{\perp}$. Since $p<k$, by the construction of $G^{(k)}_{m+k+2}$,  there exists a vertex $u$ of level $m+1$ such that $uv_i\in E(G^{(k)}_{m+k+2})$ and $\pi^{(k)}(uv_i)=0$ for each $i\in[p]$. Let  $C_1' \coloneqq {C_1} \cup \{u\}$.
    By Lemma~\ref{lemma: clique independence}, we have that $\mathbf{u}, \mathbf{v}_1, \ldots, \mathbf{v}_p$ are linearly independent.
    By $\pi^{(k)}(uv_i)=0$ for each $i\in[p]$, we have that $\mathbf{u} \perp \textbf{V}_{C_1}$. 
    Then for every $\mathbf{w} \in \textbf{V}_{C_1}$, we have $\mathbf{u} \perp \mathbf{w}$. 
    We also have $\mathbf{w} \perp \textbf{V}_{C_1}$ since $\textbf{V}_{C_1} \subseteq \textbf{V}_{C_1}^{\perp}$.  Therefore, $\mathbf{w} \perp \textbf{V}_{C_1'}$ and by the arbitrariness of $\mathbf{w}$, we have that $\textbf{V}_{C_1} \subseteq \textbf{V}_{C_1'} \cap \textbf{V}_{C_1'}^{\perp}$ which implies $\dim(\textbf{V}_{C_1'} \cap \textbf{V}_{C_1'}^{\perp})\ge p$. Hence, $C_1'$ is a bad $(p+1)$-clique,  a contradiction with the maximality of ${C_1}$, and hence, $\dim(\textbf{U})= p-1$.
    In fact, we conclude that $\textbf{U}$ is a self-orthogonal $(p-1)$-dim subspace of $\mathbb{F}_2^{2k-1}$ and each vector in $\textbf{U}$ is even. 

    Since  $\dim(\textbf{U})= p-1$, there exists $i \in \{1,2,\ldots, p\}$ such that $\mathbf{v}_i \notin \mathbf{U}$, say $i=1$.
%    there is   $\textbf{v}_i\in \{\textbf{v}_1,\textbf{v}_2,\dots,\textbf{v}_p\}$, say $i=1$, such that  $\textbf{v}_1\notin \textbf{U}$. 
   Then $\mathcal{L}(\textbf{U},\textbf{v}_1)=\textbf{V}_{C_1}$.
   If $\mathbf{v}_1$ is even, then $\mathbf{v}_1 \perp \mathcal{L}(\textbf{U},\textbf{v}_1)$, which contradicts with $\mathbf{v}_1 \notin \textbf{U}$. Thus we have that $\textbf{v}_1$ is odd. 

\begin{claim}\label{C zero neighbor is odd}
    If $u$ is a vertex in $G^{(k)}_{m+k}$ such that $uv_i\in E(G^{(k)}_{m+k})$ and $\pi^{(k)}(uv_i)=0$ for each $i\in[p]$, then $\textbf{u}$ is odd.
\end{claim}

\begin{proof}[of Claim~\ref{C zero neighbor is odd}]
Suppose $\textbf{u}$ is even. Then $\textbf{u}\perp \mathcal{L}(\textbf{u},\textbf{V}_{C_1})$. We have $\textbf{u}\in \textbf{V}_{C_1\cup\{u\}}\cap \textbf{V}_{C_1\cup\{u\}}^{\perp}$ and $\textbf{U}\subseteq \textbf{V}_{C_1\cup\{u\}}\cap \textbf{V}_{C_1\cup\{u\}}^{\perp}$. 
From Lemma~\ref{lemma: clique independence}, $\mathbf{v}_1, \ldots, \mathbf{v}_p, \mathbf{u}$ are linearly independent, so $\mathbf{u} \notin \textbf{V}_{C_1}$.
Recall that $\mathbf{U} = \mathbf{V}_{C_1} \cap \mathbf{V}_{C_1}^{\perp}$, then we have $\mathbf{u} \notin \mathbf{U}$.
Then $\dim(\textbf{V}_{C_1\cup\{u\}}\cap \textbf{V}_{C_1\cup\{u\}}^{\perp})\ge p$. Thus $ C_1\cup \{u\}$ is a bad $(p+1)$-clique, which contradicts the maximality of $C_1$. %\hfill $\square$\par
\end{proof}
\vspace{.2cm}

Fix a vertex $v_0$ of level $m+1$ such that $v_0v_i\in E(G^{(k)}_{m+k})$ and $\pi^{(k)}(v_0v_i)=0$ for each $i\in[p]$. Then $\mathbf{v}_0 \perp \mathbf{V}_{C_1}$.
By the construction of $G^{(k)}_{m+k}$, there exists a set of vertices $C_2=\{w_1,w_2,\dots, w_{k-p}\} \subseteq V(G^{(k)}_{m+k})$ satisfying the following:
\begin{enumerate}
    \item $\{v_0,v_1,\ldots,v_p,w_1,w_2,\ldots, w_{k-p}\}$ is a $(k+1)$-clique;
    \item $\pi^{(k)}(w_i w_j)=0$, for each $i,j\in[k-p], i\neq j$;
    \item and $\pi^{(k)}(v_i w_j)=\textbf{v}_i \cdot (\textbf{v}_1+\textbf{v}_0)$, for each $i\in[0,p],j\in[k-p]$.
\end{enumerate}

The existence of such $C_2$ is guaranteed because $C_1$ must be in some $k$-clique, and then we can pick $\mathbf{w}_1, \ldots, \mathbf{w}_{k-p}$ one by one according to the construction of $G^{(k)}_{m+k}$.

\begin{claim}\label{w even}
    $\textbf{w}_j$ is even for each $j\in[k-p]$.
\end{claim}

\begin{proof}[of Claim~\ref{w even}]
Suppose there is $j\in[k-p]$ such that $\textbf{w}_j$ is odd. Let  $\boldsymbol{\beta} \coloneqq \textbf{w}_j+\textbf{v}_1$. Recall that $\textbf{v}_1$ is odd.
Then 
\[
    \boldsymbol{\beta} \cdot \textbf{w}_j=\textbf{w}_j\cdot \textbf{w}_j+\textbf{v}_1\cdot \textbf{w}_j=\textbf{w}_j\cdot \textbf{w}_j+\textbf{v}_1\cdot (\textbf{v}_1+ \textbf{v}_0)=0.
\]
On the other hand, for every $v_i\in C_1$, 
\[
    \boldsymbol{\beta} \cdot \textbf{v}_i=\textbf{w}_j\cdot \textbf{v}_i+\textbf{v}_1\cdot \textbf{v}_i=\textbf{v}_i\cdot \textbf{v}_1+\textbf{v}_i\cdot \textbf{v}_0+\textbf{v}_1\cdot \textbf{v}_i=0.
\]
Hence $\boldsymbol{\beta} \perp \textbf{V}_{C_1\cup\{w_j\}}$. We have $\boldsymbol{\beta} \in \textbf{V}_{C_1\cup\{w_j\}}\cap \textbf{V}_{C_1\cup\{w_j\}}^{\perp}$ and $\textbf{U}\subseteq \textbf{V}_{C_1\cup\{w_j\}}\cap \textbf{V}_{C_1\cup\{w_j\}}^{\perp}$. From Lemma~\ref{lemma: clique independence} and $\dim(\textbf{U})= p-1$, we know $\textbf{w}_j \notin \textbf{U}$. Then $\dim(\textbf{V}_{C_1\cup\{w_j\}}\cap \textbf{V}_{C_1\cup\{w_j\}}^{\perp})\ge p$ which implies  $ C_1\cup \{w_j\}$ is a bad $(p+1)$-clique, a contradiction with the maximality of $C_1$.
\end{proof}

\vspace{.2cm}

\begin{table}
    \centering
    \begin{tabular}{|c|c|c|c|c|c|c|c|}
    \hline
         & $\textbf{v}_0$ & $\textbf{v}_1$  & $\textbf{v}_i$ & $\textbf{w}_{j_1}$ & $\textbf{w}_{j_2}$ & $\boldsymbol{\beta}_{j_1}$ &  $\boldsymbol{\beta}_{j_2}$\\
    \hline
        $\textbf{v}_0$ & 1 & 0 & 0 & 1 & 1 & 0 & 0\\
    \hline
        $\textbf{v}_1$ & 0 & 1 & $\textbf{v}_1\cdot \textbf{v}_i$ & 1 & 1 & 0 & 0\\
    \hline
        $\textbf{v}_i$ & 0 & $\textbf{v}_1\cdot \textbf{v}_i$ & $\textbf{v}_i\cdot \textbf{v}_i$ & $\textbf{v}_1\cdot \textbf{v}_i$ & $\textbf{v}_1\cdot \textbf{v}_i$ & 0 & 0\\
    \hline
        $\textbf{w}_{j_1}$ & 1 & 1 & $\textbf{v}_1\cdot \textbf{v}_i$ & 0 & 0 & 0 & 0\\
    \hline
        $\textbf{w}_{j_2}$ & 1 & 1 & $\textbf{v}_1\cdot \textbf{v}_i$ & 0 & 0 & 0 & 0\\
    \hline
        $\boldsymbol{\beta}_{j_1}$ & 0 & 0 & 0 & 0 & 0 & 0 & 0\\
    \hline
        $\boldsymbol{\beta}_{j_2}$ & 0 & 0 & 0 & 0 & 0 & 0 & 0\\
    \hline
    \end{tabular}
    \caption{The inner products for $i\in[p]$ and $j_1,j_2\in[k-p]$}\label{tab}
\end{table}

Now we complete the proof of Lemma~\ref{lemma: 0 sol independence}. For each $j\in[k-p]$, let $\boldsymbol{\beta}_j=\textbf{v}_0+\textbf{v}_1+\textbf{w}_j$. By Claim~\ref{C zero neighbor is odd},  $\textbf{v}_{0}$ is odd. By Claim~\ref{w even}, $\textbf{w}_j$ is even for each $j\in[k-p]$.
 It is not difficult to show that  $\boldsymbol{\beta}_j\perp \textbf{v}_{0}$, $\boldsymbol{\beta}_j\perp \textbf{V}_{C_1}$ and $\boldsymbol{\beta}_j\perp \textbf{V}_{C_2}$. Check Table~\ref{tab} for the inner products between the vectors that we are working with.
Let $C_3=C_1\cup C_2$ and $\mathbf{W}=\textbf{V}_{C_3}=\textbf{V}_{C_1}+\textbf{V}_{C_2}$. Then  $\dim(\mathbf{W})=k$ from Lemma~\ref{lemma: clique independence}. Since $\mathbf{W} \subseteq  \mathbb{F}_2^{2k-1}$, we have  $\dim(\mathbf{W} ^{\perp})=k-1$. Let $\mathbf{W} ^{\prime}=\mathcal{L}(\textbf{U},\boldsymbol{\beta}_1,\dots,\boldsymbol{\beta}_{k-p})$. Since $\textbf{U}\perp \mathbf{W}$ and $\boldsymbol{\beta}_j\perp \mathbf{W}$ for each $j\in[k-p]$, we have  $\mathbf{W}^{\prime}\subseteq \mathbf{W}^{\perp}$. Note that $\textbf{V}_{C_1},\textbf{V}_{C_2}\subseteq \mathcal{L}(\textbf{v}_0,\textbf{v}_1,\mathbf{W}^{\prime})$. We have $\mathbf{W} \subseteq \mathcal{L}(\textbf{v}_0,\textbf{v}_1,W^{\prime})$ which implies  $\dim(\mathbf{W}^{\prime})\ge k-2$. If $\textbf{v}_0\notin \mathbf{W}$, then  $\dim(\mathbf{W}^{\prime})\ge k-1$ which implies $\mathbf{W}^{\perp}=\mathbf{W}^{\prime}$. Since $\textbf{v}_0\perp \mathbf{W}^{\prime}$, we have $\textbf{v}_0\perp \mathbf{W}^{\perp}$, a contradiction with $\textbf{v}_0\notin \mathbf{W}$. Hence $\textbf{v}_0\in \mathbf{W}$.
By Lemma~\ref{lemma: independent or sum}, we have $ \textbf{v}_0+\cdots+\textbf{v}_p+\textbf{w}_1+\cdots+\textbf{w}_{k-p}=0 $.

Since $\textbf{v}_0\in \mathbf{W}$, we have $\boldsymbol{\beta}_j\in \mathbf{W}$ for each $j$ which implies $\mathbf{W}^{\prime}\subseteq \mathbf{W}$. So  $\mathbf{W}^{\prime}\subseteq \mathbf{W}\cap \mathbf{W}^{\perp}$.
If  $\dim(\mathbf{W}^{\prime})\ge k-1$, then $C_3$ is a bad $k$-clique,  a contradiction with the maximality of $C_1$. Hence $\dim(\mathbf{W}^{\prime})=k-2$. Let $\boldsymbol{\alpha} \in \mathbf{W}^{\perp}\backslash \mathbf{W}^{\prime}$ such that $\mathbf{W}^{\perp}=\mathcal{L}(\mathbf{W}^{\prime},\boldsymbol{\alpha})$. By the construction of $G^{(k)}_{m+k}$, there exists  a vertex $x$ connecting to all vertices of $C_3$ such that $\pi^{(k)}(xy) = 0$ for each $y \in C_3$. Then $\textbf{x}\in \mathbf{W}^{\perp}$.
From Claim~\ref{C zero neighbor is odd}, $\textbf{x}$ is odd. 
Since $\textbf{U}$ is a self-orthogonal subspace and $\boldsymbol{\beta}_i\bot \boldsymbol{\beta}_j$ for every $i,j\in [k-p]$ (see Table 1), we have that
the vectors in $\mathbf{W}^{\prime}$ are all even.  By $\textbf{x}\in \mathbf{W}^{\perp}=\mathcal{L}(\mathbf{W}^{\prime},\boldsymbol{\alpha})$ and $\textbf{x}$ being odd, we have $\boldsymbol{\alpha}$ is odd. 
Let $C^*=\{v_0,\ldots,v_p,w_1,\ldots,w_{k-p-1}\}$ (if $p=k-1$, then let $C^*=\{v_0,\ldots,v_p\}$). Then there is 
 $x^*$ connecting to all vertices of $C^*$ such that $\pi^{(k)}(x^* y) =\textbf{v}_0\cdot \textbf{y}$ for each $y\in C^*$. 
 
 Recall that $\mathbf{W} = \mathcal{L}(\mathbf{v}_1, \ldots, \mathbf{v}_p, \mathbf{w}_1, \ldots, \mathbf{w}_{k-p})$ and $\mathbf{V}_{C^*} = \mathcal{L}(\mathbf{v}_0, \ldots, \mathbf{v}_p, \mathbf{w}_1, \ldots, \mathbf{w}_{k-p-1})$. 
 Since $ \textbf{v}_0+\cdots+\textbf{v}_p+\textbf{w}_1+\cdots+\textbf{w}_{k-p}=0 $, we have that $\mathbf{W}=\textbf{V}_{C^*}$.
  Then $\textbf{x}^*\in \textbf{v}_0 + \mathbf{W}^{\perp}$. Since $\pi^{(k)}(x^* v_i) =\textbf{v}_0\cdot \textbf{v}_i=0$ for each $i\in [k]$,
  from Claim~\ref{C zero neighbor is odd}, $\textbf{x}^*$ is odd. Note that $\textbf{x}^*\in \textbf{v}_0+\mathcal{L}(\boldsymbol{\alpha},\mathbf{W}^{\prime})$. Since all vectors in $\mathbf{W}^{\prime}$ are  even and $\textbf{v}_0,\boldsymbol{\alpha}$ are odd, we have $\textbf{x}^*\in \textbf{v}_0+\mathbf{W}^{\prime}\subseteq \mathbf{W}$.
  From Lemma~\ref{lemma: independent or sum}, $\textbf{x}^*= \textbf{v}_0+\cdots+\textbf{v}_p+\textbf{w}_1+\cdots+\textbf{w}_{k-p-1}=\textbf{w}_{k-p}$. Since $\textbf{x}^{*}\cdot \mathbf{v}_1=0\neq 1=\textbf{w}_{k-p}\cdot \mathbf{v}_1$, we derive a contradiction.
\end{proof}

\end{document}
